\documentclass[11pt,letterpaper]{article}
\usepackage[T1]{fontenc}
\usepackage[utf8]{inputenc}
\usepackage{lmodern}
\usepackage[margin=1in,headheight=15pt]{geometry}
\usepackage{amsmath,amssymb,amsthm,mathtools,mathrsfs}
\usepackage{microtype,booktabs,longtable,array,enumitem}
\usepackage[dvipsnames]{xcolor}
\usepackage{fancyhdr,aliascnt}
\usepackage[colorlinks=true,linkcolor=MidnightBlue,citecolor=MidnightBlue,urlcolor=MidnightBlue]{hyperref}
\usepackage[nameinlink,noabbrev]{cleveref}
\hypersetup{pdftitle={Global Support Obstructions in Surcomplex Analysis},pdfauthor={Research manuscript prepared with ChatGPT},pdfsubject={Moving divisors, Hahn Mittag-Leffler theory, and a Picard-group dichotomy}}
\pagestyle{fancy}
\fancyhf{}
\fancyhead[L]{\small Global support in surcomplex analysis}
\fancyhead[R]{\small Divisors and line bundles}
\fancyfoot[C]{\thepage}
\setlength{\parskip}{0.3em}
\setlength{\parindent}{1.25em}
\setlist{nosep,leftmargin=2em}
\newtheorem{theorem}{Theorem}[section]
\newaliascnt{lemma}{theorem}
\newtheorem{lemma}[lemma]{Lemma}
\aliascntresetthe{lemma}
\newaliascnt{proposition}{theorem}
\newtheorem{proposition}[proposition]{Proposition}
\aliascntresetthe{proposition}
\newaliascnt{corollary}{theorem}
\newtheorem{corollary}[corollary]{Corollary}
\aliascntresetthe{corollary}
\theoremstyle{definition}
\newaliascnt{definition}{theorem}
\newtheorem{definition}[definition]{Definition}
\aliascntresetthe{definition}
\newaliascnt{example}{theorem}
\newtheorem{example}[example]{Example}
\aliascntresetthe{example}
\theoremstyle{remark}
\newaliascnt{remark}{theorem}
\newtheorem{remark}[remark]{Remark}
\aliascntresetthe{remark}
\newcommand{\No}{\mathbf{No}}
\newcommand{\SC}{\mathbf{SC}}
\newcommand{\C}{\mathbb C}
\newcommand{\R}{\mathbb R}
\newcommand{\N}{\mathbb N}
\newcommand{\Z}{\mathbb Z}
\newcommand{\Q}{\mathbb Q}
\newcommand{\OO}{\mathcal O}
\newcommand{\HH}{\mathscr H_\Gamma}
\newcommand{\HP}{\mathscr H_{\Gamma,+}}
\newcommand{\HI}{\mathscr H_{\Gamma,\ge0}}
\newcommand{\KK}{K_\Gamma}
\newcommand{\mm}{\mathfrak m_\Gamma}
\newcommand{\sh}{^{\sharp}}
\newcommand{\st}{\operatorname{st}}
\newcommand{\supp}{\operatorname{supp}}
\newcommand{\ord}{\operatorname{ord}}
\newcommand{\Pic}{\operatorname{Pic}}
\newcommand{\Div}{\operatorname{Div}}
\newcommand{\ExpH}{\operatorname{Exp}_{H}}
\newcommand{\LogH}{\operatorname{Log}_{H}}
\newcommand{\id}{\operatorname{id}}
\newcommand{\PP}{\operatorname{pp}}
\newcommand{\abs}[1]{\left|#1\right|}
\newcommand{\mgen}[1]{\langle #1\rangle_{+}}
\newcolumntype{L}[1]{>{\raggedright\arraybackslash}p{#1}}
\numberwithin{equation}{section}
\begin{document}
\begin{titlepage}
\centering
\vspace*{0.7cm}
{\Large\scshape A research continuation\par}
\vspace{0.65cm}
{\Huge\bfseries Global Support Obstructions\par}
\vspace{0.25cm}
{\LARGE\bfseries in Surcomplex Analysis\par}
\vspace{0.65cm}
{\Large Moving Divisors, Mittag--Leffler Theory,\par and a Picard-Group Dichotomy\par}
\vspace{0.8cm}
{\large September 21, 2026\par}
\vspace{0.6cm}
\begin{minipage}{0.93\textwidth}
\begin{abstract}
We study the sheaf $\HH$ of well-ordered Hahn series with ordinary holomorphic coefficient functions, in a fixed set-sized surreal value group $\Gamma$. The question is global: when do infinitely many individually legitimate infinitesimal analytic data admit one coherent realization? We prove that a Hahn Mittag--Leffler problem on a discrete subset of the complex plane is solvable exactly when the union of its principal-part supports is well ordered. Combining this criterion with positive-support logarithms gives a sharp Weierstrass theorem for moving surcomplex divisors: prescribed finite clusters over a discrete ordinary divisor are globally realizable exactly when the supports of their monic cluster-polynomial coefficients have a well-ordered union.

The obstruction detects symmetric cancellation. Simple zeros at $j+\omega^{-1/j}$ cannot be realized alone by a coherent entire function, whereas the complete clusters $(z-j)^j=\omega^{-1}$ can. Thus a globally principal effective divisor can have a nonprincipal effective sub-divisor. We also show that every ordinary-base stalk of $\HH$ is a principal ideal domain, although generally not a local ring.

Finally, positive Hahn exponentiation identifies the nonclassical line-bundle obstruction with $H^1(\C,\HP)$. The resulting Picard group vanishes exactly when $\Gamma$ is trivial or cyclic. For every noncyclic value group it contains a complex vector subspace of dimension $2^{\aleph_0}$; for countable noncyclic groups this is its exact dimension. These statements concern the ordinary analytic base and strong Hahn summability, not the fine topology of the full surreal class.
\end{abstract}
\end{minipage}
\vfill
\begin{minipage}{0.93\textwidth}\small
\textbf{Status.} The principal global criteria and the Picard-group dichotomy are proposed new results in the precise framework of the supplied manuscript. Full proofs are given. The literature check did not locate these formulations, but does not certify priority. No established named open conjecture, peer review, or machine-checked proof is claimed.
\end{minipage}
\end{titlepage}
\tableofcontents
\newpage

\section{The question: from local roots to global compatibility}\label{sec:intro}
The supplied manuscript \cite{Source} develops surcomplex analysis by separating fine-local power-series calculus from Hahn-coherent holomorphy over an ordinary complex domain. Its one-variable preparation theorem associates a finite monic polynomial to every infinitesimal zero cluster. Its root-lifting theorem counts the roots in a monad exactly. Those are the starting points here, not new claims of this article.

A remaining question is not answered by local preparation. Suppose an ordinary discrete set has infinitely many points, and a finite infinitesimal root configuration has been specified above each point. Can one coherent entire function realize all these configurations simultaneously? The classical Weierstrass theorem might suggest that ordinary discreteness is sufficient. It is not. A condition on the union of the \emph{Hahn supports} is necessary, and it turns out to be sufficient as well.

The same issue appears in the additive Cousin, or Mittag--Leffler, problem. Solving each complex coefficient separately is easy. Assembling those solutions into one Hahn series may be impossible. Exponentiating the positive-support version of this obstruction produces nontrivial line bundles even though the ordinary base is the contractible Stein manifold $\C$.

\subsection{The main results in concrete form}
Fix an ordered additive subgroup $\Gamma\subseteq\No$ that is a set, and write $t^\gamma=\omega^{-\gamma}$. A section of $\HH$ on a connected ordinary domain $U$ is
\[
 F(z)=\sum_{\gamma\in S} f_\gamma(z)t^\gamma,
 \qquad f_\gamma\in\OO(U),\qquad S\subseteq\Gamma\ \text{well ordered}.
\]
Here the variable $z$ belongs to the ordinary analytic base at the coefficient level; evaluation on finite surcomplex arguments uses Taylor expansion. The precise definitions are in \cref{sec:setup}.

The principal conclusions are as follows.
\begin{enumerate}
\item \textbf{Sharp principal-part compatibility.} A discrete family of Hahn principal parts is the family of principal parts of a single section on the punctured plane if and only if all the principal-part exponents together form a well-ordered set (\cref{thm:ML}).
\item \textbf{Sharp moving-divisor realization.} Write the prescribed cluster over $c_j$ as the zero set of
\[
 P_j(u)=u^{m_j}+\sum_{r<m_j}a_{j,r}u^r,
 \qquad a_{j,r}\in\mm.
\]
One coherent entire function has exactly these clusters, with their prescribed multiplicities, if and only if $\bigcup_{j,r}\supp(a_{j,r})$ is well ordered (\cref{thm:divisor}). It is the support of the \emph{symmetric coefficients}, not the supports of separately chosen roots, that matters.
\item \textbf{A line-bundle dichotomy.} With the usual tensor-product operation on invertible sheaves,
\[
 \Pic(\C,\HH)=0
 \quad\Longleftrightarrow\quad
 \Gamma=\{0\}\ \text{or}\ \Gamma=\Z\delta\quad(\delta>0).
\]
For countable noncyclic $\Gamma$, this Picard group has a natural complex vector-space structure of dimension $2^{\aleph_0}$ (\cref{thm:dichotomy,thm:dimension}). The interpretation of invertible sheaves is justified by the principal-ideal-domain theorem for stalks in \cref{sec:stalks}.
\end{enumerate}
These are global existence and obstruction theorems, not just formal analogues of differentiation identities. They distinguish finite from infinite collections of local data and cyclic from divisible or higher-rank value groups.

\subsection{Relation to the attachment and to published work}\label{sub:literature}
The terminology $U\sh$, strong summability, and Hahn coherence follows \cite{Source}. We preserve the source's distinction between the ordinary base and the fine topology. We work in one fixed workspace to make sheaves and cohomology set-valued. The resulting theorems apply in particular to $\Gamma=\Q$, whose Hahn field embeds in $\SC=\No[i]$.

Normal forms, infinitesimal power-series evaluation, and the underlying surreal field theory have established predecessors, including van den Dries--Ehrlich \cite{vdDE}; the algebraic closure of divisible Hahn fields is covered by Poonen \cite{Poonen}. Berarducci--Mantova's surreal function theory \cite{BM} and Costin--Ehrlich's integration theory \cite{CE} address different extension and summation questions. We do not identify their function classes with the sheaf studied here. Published surcomplex exponentiation also exists: Ehrlich--Kaplan construct canonical extensions for their trigonometric-exponential fields \cite[\S11]{EK}. No novelty claim about merely defining a surcomplex exponential is made here.

There is another important terminological distinction. M\"uller--Strohmaier's Hahn-holomorphic and Hahn-meromorphic functions have normally convergent generalized expansions near a point on a logarithmic cover \cite{MS}. In our notation, the Hahn parameter is an independent surreal scale, and the coefficients are holomorphic functions on an ordinary base. There is no ordinary convergence requirement across the Hahn exponents. Similar terminology does not make the objects or the global support conditions identical.

The classical analytic inputs are the Weierstrass product theorem, Mittag--Leffler's theorem, entire interpolation, and the vanishing of $H^1(\C,\OO)$ and $\Pic(\C,\OO)$; see \cite{Ahlfors,Demailly}. These inputs are explicitly separated from the Hahn arguments. The present proofs do not apply Cartan's theorem to $\HH$ as though it were a coherent $\OO$-module.

\subsection{Novelty and scope of the contribution}
The requested alternative of proving new theorems, rather than solving a named published conjecture, is the route taken here. The principal claims of possible novelty are the exact support criteria, their symmetric-cluster consequences, the explicit obstruction subgroup, and the value-group dichotomy. The local preparation theorem is inherited and reproved with its uniform support bound. The stalk theorem and near-identity interpolation theorem are additional deductions needed to clarify the structure.

The literature search conducted on September 21, 2026 included combinations of ``surcomplex'', ``Hahn holomorphic'', ``Mittag--Leffler'', ``moving divisors'', ``Picard group'', and ``holomorphic formal Laurent series'', together with the primary sources cited above. No source located in that search stated the main results for this sheaf. This is a qualified statement about the search, not evidence of an exhaustive bibliographic classification. The mathematical assertions stand or fall by the proofs below.

\section{The coefficient sheaf and its support algebra}\label{sec:setup}
\subsection{Fixed Hahn fields and infinitesimals}
Let $\Gamma$ be a set-sized ordered additive subgroup of $\No$. Its complex Hahn field is
\[
 \KK=\C((t^\Gamma))
 =\left\{\sum_{\gamma\in S}a_\gamma t^\gamma:
 S\subseteq\Gamma\text{ well ordered},\ a_\gamma\in\C\right\}.
\]
Zero coefficients are omitted. Conway normal forms embed this field in $\SC$ by $t^\gamma=\omega^{-\gamma}$. For $a\ne0$, set $v(a)=\min\supp(a)$, and let $v(0)=+\infty$. We use
\[
 \mm=\{a\in\KK:v(a)>0\},
 \qquad \OO_\Gamma=\{a\in\KK:v(a)\ge0\}.
\]
The zero element is included in $\mm$. Elements of $\mm$ are infinitesimal in the surreal modulus, while $\OO_\Gamma$ consists of the finite elements of this workspace. The coefficient of $t^0$ gives the standard part of a finite element.

We use algebraic closedness of $\SC$, and occasionally of $\KK$ when $\Gamma$ is divisible \cite{Poonen}. No divisibility assumption is needed for the support, preparation, Mittag--Leffler, or Picard theorems. Roots of polynomials can always be discussed in $\SC$, or in a divisible Hahn enlargement.

\begin{lemma}[Support calculus]\label{lem:support}
If $S,T$ are well-ordered subsets of an ordered abelian group, then $S+T$ is well ordered and every element has finitely many representations $s+t$. If $E$ is well ordered and $E\subseteq\Gamma_{>0}$, then
\[
 E^*:=\{e_1+\cdots+e_n:n\ge0,\ e_i\in E\}
\]
is well ordered, and an exponent has only finitely many representations by nondecreasing finite words in $E$. Consequently, a formal power series in finitely many positive-support Hahn elements is strongly summable.
\end{lemma}
This is the Hahn--Neumann lemma \cite{Neumann,Poonen}. The finiteness assertion refers to words in the \emph{set of exponents}, not to an infinite family of spatial points carrying the same exponent. Finite permutations account for ordered products. We repeatedly use both conclusions: a well-ordered support alone is not enough to justify an arbitrary rearrangement if infinitely many summands contribute at one exponent.

A useful illustration of the non-Archimedean issue is that $\gamma>0$ need not make $n\gamma$ eventually exceed every element of $\Gamma$. Our arguments do not use such an Archimedean property. They use finite decomposition in $E^*$.

\subsection{Sections and ordinary-base gluing}
\begin{definition}[Hahn-coherent coefficient sheaf]
For connected nonempty ordinary open $U\subseteq\C$, define
\begin{equation}\label{eq:Hdef}
 \HH(U)=\OO(U)((t^\Gamma))
 =\left\{\sum_{\gamma\in S}f_\gamma t^\gamma:
 f_\gamma\in\OO(U),\quad S\subseteq\Gamma\text{ well ordered}\right\}.
\end{equation}
For disconnected $U$, take the product over its connected components. Define $\HP$ by requiring strictly positive exponents, and $\HI$ by requiring nonnegative exponents, componentwise. These are respectively an additive sheaf with multiplication and a sheaf of unital subrings of $\HH$.
\end{definition}
The support of a datum is the set of exponents whose coefficient function is \emph{not identically zero}, not the support of its value at a particular point. Derivatives differentiate coefficient functions and leave $t^\gamma$ constant.

\begin{proposition}[Support-preserving sheaf property]\label{prop:sheaf}
The preceding prescriptions are sheaves. On a connected ordinary domain, compatible local sections glue with a single well-ordered support.
\end{proposition}
\begin{proof}
Refine a cover to connected opens. On an overlap, equality of Hahn data implies equality of each coefficient function. A holomorphic coefficient nonzero on a connected open cannot vanish identically on a nonempty open subset. Therefore two sections on intersecting connected opens have the same exact support if they agree on the overlap. The intersection graph of a cover of a connected space by connected opens is connected: otherwise its graph components would give a separation of the space. Thus all nonempty local descriptions share one support. Glue each holomorphic coefficient and retain that support. Uniqueness is coefficientwise. Positivity and nonnegativity are preserved. This is the gluing mechanism already used in \cite[\S5]{Source}.
\end{proof}
This proof explains why arbitrary unions of supports are \emph{not} allowed merely because a problem is local. Equality on overlaps provides extra rigidity; an additive or multiplicative cocycle does not itself give that equality.

\subsection{Evaluation on surcomplex monads}
For ordinary $U$, write
\[
 U\sh=\{c+\varepsilon:c\in U,\ \varepsilon\in\SC\text{ infinitesimal}\}.
\]
The evaluation of \eqref{eq:Hdef} is
\begin{equation}\label{eq:eval}
 F\sh(c+\varepsilon)
 =\sum_{\gamma\in S}\sum_{n\ge0}
 t^\gamma\frac{f_\gamma^{(n)}(c)}{n!}\varepsilon^n.
\end{equation}
The support lies in $S+(\supp\varepsilon)^*$, and \cref{lem:support} proves strong summability with finite contributions at each exponent. Evaluation is injective, since evaluation at all ordinary points recovers every coefficient function. It respects addition, multiplication, and coefficientwise differentiation; ordinary Taylor identities and the same support estimate prove these assertions. These are the evaluation results of \cite[\S5]{Source}.

On the workspace thickening, where $\varepsilon\in\mm$, the output lies in $\KK$. Allowing infinitesimals outside the workspace changes the evaluation domain, not the coefficient sheaf or its cohomology. We count roots in the full monad in $\SC$, unless explicitly restricting the field.

For $F=f_0+E$ with $E\in\HP(U)$ and $f_0(U)\subseteq V$, composition with $H=\sum h_\delta t^\delta\in\HH(V)$ is
\begin{equation}\label{eq:comp}
 H\circ F=\sum_{\delta,n\ge0}t^\delta
 \frac{h_\delta^{(n)}\circ f_0}{n!}E^n.
\end{equation}
Its support is contained in $\supp H+(\supp E)^*$; evaluation gives actual composition. The analogous formula at an infinitesimal argument is used throughout. There is no claim that ordinary partial sums converge in the full fine topology.

\section{Units, positive logarithms, and the integral subring}\label{sec:units}
\begin{proposition}[Exact unit criterion]\label{prop:units}
A nonzero section $F\in\HH(U)$ on connected $U$ is a unit if and only if its leading coefficient function is nowhere zero. Every such unit has a unique factorization
\begin{equation}\label{eq:unitfactor}
 F=t^\lambda h(1+E),\qquad
 \lambda\in\Gamma,\quad h\in\OO(U)^\times,\quad E\in\HP(U).
\end{equation}
\end{proposition}
\begin{proof}
The valuation of a product is the sum of the leading exponents because $\OO(U)$ is a domain. If $FG=1$, the leading coefficient functions multiply to $1$, so the leading coefficient of $F$ has no zeros. Conversely, division by that leading coefficient and by the leading monomial gives $1+E$ with positive support. The series $\sum_{n\ge0}(-E)^n$ is strongly summable by \cref{lem:support} and is its inverse. Uniqueness follows by comparing the leading exponent and coefficient.
\end{proof}

\begin{proposition}[Positive-support exponential and logarithm]\label{prop:explog}
The formulas
\begin{equation}\label{eq:explog}
 \ExpH A=\sum_{n\ge0}\frac{A^n}{n!},\qquad
 \LogH(1+E)=\sum_{n\ge1}\frac{(-1)^{n+1}E^n}{n}
\end{equation}
define mutually inverse isomorphisms of abelian sheaves
\[
 (\HP,+)\simeq(1+\HP,\cdot).
\]
Their supports are contained in $(\supp A)^*$ and $(\supp E)^*\setminus\{0\}$, respectively, with the constant term of the exponential treated separately. In particular,
\begin{equation}\label{eq:unitsplit}
 \HH^\times\simeq\underline\Gamma\times\OO^\times\times\HP,
 \qquad (\lambda,h,A)\longmapsto t^\lambda h\ExpH A.
\end{equation}
Here $\underline\Gamma$ denotes the locally constant sheaf with values in the additive group $\Gamma$.
\end{proposition}
\begin{proof}
The supports of all displayed terms are controlled by a positive monoid, with finite word decompositions. Hence the formal characteristic-zero identities for logarithm and exponential can be evaluated and regrouped. They give the inverse identities and $\ExpH(A+B)=\ExpH(A)\ExpH(B)$. Apply \cref{prop:units} componentwise for \eqref{eq:unitsplit}. All constructions commute with restriction.
\end{proof}
There is no selection of a global logarithm of an arbitrary ordinary holomorphic unit in this decomposition. That factor remains in $\OO^\times$. Only the positive-support factor is logarithmized canonically.

\begin{remark}[Reduction is not defined on the full coefficient field]\label{rem:reduction}
Taking the coefficient at exponent zero is a ring homomorphism $\HI\to\OO$, with kernel $\HP$. It is not a ring homomorphism $\HH\to\OO$: for $\delta>0$, both $t^\delta$ and $t^{-\delta}$ would map to zero while their product is $1$. Later, a line bundle with transitions in $1+\HP$ has an integral model whose reduction is the trivial ordinary line bundle. This statement does not assert a reduction map on the full ring $\HH$.
\end{remark}

\section{Uniform local preparation and the ordinary-base stalks}\label{sec:stalks}
We record the local result from the supplied manuscript with the support bound required for an infinite family of centers. This avoids treating local preparation as a black box whose supports might depend uncontrollably on the center.

\subsection{Preparation with one controlling monoid}
For $h$ holomorphic on a disk about zero, let
\[
 R_mh=\sum_{r=0}^{m-1}\frac{h^{(r)}(0)}{r!}u^r,
 \qquad D_mh=\frac{h-R_mh}{u^m}.
\]
Both are holomorphic, and $h=R_mh+u^mD_mh$.

\begin{theorem}[Uniform-support local preparation]\label{thm:preparation}
Let $F\in\HH(U)$ be nonzero, with leading exponent $\lambda$ and leading coefficient $f_\lambda$. Put
\[
 T=(\supp F-\lambda)\setminus\{0\}\subseteq\Gamma_{>0},
 \qquad M=T^*.
\]
If $f_\lambda$ has a zero of order $m\ge1$ at $c\in U$, there is an ordinary disk $D$ about zero on which
\begin{equation}\label{eq:prep}
 F(c+u)=t^\lambda q_0(u)P_c(u)Q_c(u),
\end{equation}
where $q_0(u)=f_\lambda(c+u)/u^m$ is an ordinary holomorphic unit,
\[
 P_c(u)=u^m+\sum_{r<m}a_{c,r}u^r,
 \qquad a_{c,r}\in\mm,
 \qquad Q_c\in1+\HP(D).
\]
Every support $\supp(a_{c,r})$ and $\supp(Q_c-1)$ lies in $M\setminus\{0\}$. The monoid $M$ depends on $F$, not on $c$ or $m$. The monic polynomial $P_c$ and the normalized factor $Q_c$ are unique.
\end{theorem}
\begin{proof}
Choose $D$ so $q_0$ has no zeros. Dividing by $t^\lambda q_0$ gives $u^m+E$, where the support of $E$ is contained in $T$. Seek $P_c=u^m+A$ and $Q_c=1+B$ with $\deg_u A<m$ and positive supports in $M$. The equation is
\begin{equation}\label{eq:prep-rec}
 E=A+u^m B+AB.
\end{equation}
At $\nu\in M\setminus\{0\}$, after all smaller exponents have been treated, set
\[
 h_\nu=E_\nu-
 \sum_{\substack{\alpha+\beta=\nu\\\alpha,\beta>0}}A_\alpha B_\beta,
 \qquad A_\nu=R_mh_\nu,\qquad B_\nu=D_mh_\nu.
\]
The sum is finite and all its indices are smaller than $\nu$. Well-founded recursion gives holomorphic coefficient functions on the same disk. The fixed well-ordered set $M$ contains the resulting supports; the finitely many coefficients of $A$ have positive supports and belong to $\mm$. Equation \eqref{eq:prep-rec} holds at every exponent.

For uniqueness, compare two normalized pairs and choose the least exponent at which $A$ or $B$ differs. Products involving a positive exponent and a difference cannot contribute at that least exponent. Hence the difference satisfies $\Delta A_\nu+u^m\Delta B_\nu=0$. Taking $R_m$ and $D_m$ makes both terms zero, a contradiction. A finite union of the two candidate supports is well ordered, so this comparison does not assume a common monoid in advance. This is the preparation argument of \cite[\S7]{Source}, with its center-independent support bound made explicit.
\end{proof}

\begin{corollary}[All roots in a monad]\label{cor:roots}
The zeros of $F\sh$ in the monad $c+\mathfrak m_{\SC}$ are exactly $c$ plus the roots of $P_c$, with polynomial multiplicity. Their total multiplicity is $\ord_c f_\lambda$. If $f_\lambda(c)\ne0$, there are no zeros in that monad. When $\Gamma$ is divisible, all these roots lie in $c+\mm$.
\end{corollary}
\begin{proof}
The evaluated factors $q_0$ and $Q_c$ in \eqref{eq:prep} are nonzero throughout a sufficiently small ordinary thickened disk: their relevant leading ordinary values are nonzero and $Q_c$ is $1$ plus an infinitesimal. Every root $\xi$ of a monic polynomial $u^m+\sum_{r<m}a_ru^r$ with infinitesimal lower coefficients is infinitesimal. Otherwise $1/\xi$ would be finite and
\[
 0=P_c(\xi)/\xi^m=1+\sum_{r<m}a_r\xi^{r-m}
\]
would be $1$ plus an infinitesimal. Algebraic closedness of $\SC$ now gives exactly $m$ roots with multiplicity. If $f_\lambda(c)\ne0$, evaluation has a nonzero leading term at every point of its monad. The divisible-field assertion follows from algebraic closedness of $\KK$ \cite[Corollary 4]{Poonen}.
\end{proof}
In particular, any two monic infinitesimal cluster polynomials with the same surcomplex roots and multiplicities are equal. This elementary uniqueness is important when passing from a divisor to its coefficients.

\subsection{A principal ideal domain that is not local}
Let $R_c=(\HH)_c$ be the stalk at the \emph{ordinary} base point $c$. This stalk is formed by shrinking ordinary neighborhoods. Such shrinking never separates two surcomplex points in the same monad.

\begin{theorem}[Principal-ideal-domain stalks]\label{thm:PID}
For every $c\in\C$, the domain $R_c$ is a principal ideal domain. Every nonzero element is associated to a unique monic polynomial in $u=z-c$ whose lower coefficients are infinitesimal. If $\Gamma\ne\{0\}$, $R_c$ is not a local ring. If $\Gamma$ is divisible, its maximal ideals are precisely
\begin{equation}\label{eq:maxideals}
 (u-\varepsilon),\qquad \varepsilon\in\mm,
\end{equation}
and their residue fields are $\KK$.
\end{theorem}
\begin{proof}
The ring is a domain because nonzero coefficient functions remain nonzero on smaller connected neighborhoods. Preparation shows that every nonzero germ is a unit times a distinguished monic polynomial; a unit is represented by the polynomial $1$. Uniqueness follows from \cref{cor:roots}, or from normalized preparation after multiplying by a unit. Define $d(F)$ to be the degree of this polynomial. It is a nonnegative integer and vanishes exactly on units.

First consider two nonzero germs $F,G$, associated to polynomials $P,Q\in\KK[u]$. Let $D$ be their monic polynomial greatest common divisor. The Euclidean algorithm in $\KK[u]$ gives $A P+B Q=D$, and $D$ divides both $P$ and $Q$. Hence $(F,G)=(D)$ in $R_c$. Moreover, $D$ is distinguished: in an algebraically closed Hahn enlargement all its roots are roots of $P$, hence infinitesimal, and its lower coefficients are finite elementary symmetric sums of such roots. Thus $d(D)=\deg D$.

Now let $I$ be any nonzero ideal of $R_c$. Choose $F\in I\setminus\{0\}$ for which $d(F)$ is minimal. For every nonzero $G\in I$, the preceding greatest common divisor belongs to $I$ and has degree at most $d(F)$. Minimality makes its degree equal to $d(F)$. The monic polynomial associated to $F$ therefore divides the one associated to $G$, so $F$ divides $G$ in $R_c$. Consequently $I=(F)$. This treats arbitrary ideals, not only finitely generated ones.

Choose $\delta\in\Gamma_{>0}$ when $\Gamma\ne0$. Both $u$ and $u-t^\delta$ are nonunits, since their leading ordinary coefficient is $u$. Their difference is the unit $t^\delta$. The nonunits are therefore not an ideal, and $R_c$ is not local.

If $\Gamma$ is divisible, every distinguished polynomial factors into linear factors $u-\varepsilon$ with $\varepsilon\in\mm$. Each such factor is a nonunit and is irreducible, by degree. A nonzero maximal ideal in the principal ideal domain must be generated by one of these factors. Conversely, evaluation at $c+\varepsilon$ gives a surjection $R_c\to\KK$, because constants are available. Preparation and polynomial division show that its kernel is $(u-\varepsilon)$. This proves both \eqref{eq:maxideals} and the residue-field assertion.
\end{proof}
The case $\Gamma=0$ reduces to the familiar local ring of ordinary one-variable holomorphic germs. For nonzero $\Gamma$, the extra maximal ideals record the infinitesimally separated points that the ordinary base has collapsed into one point. There is no contradiction with the fine-germ ring at a \emph{single} surcomplex point considered in \cite{Source}: the two stalk constructions use different neighborhoods.

\begin{corollary}[Meaning of the Picard group here]\label{cor:Picmeaning}
For the ringed space $(\C,\HH)$, a tensor-invertible sheaf of modules is locally free of rank one. Thus its Picard group is classified by unit-valued transition cocycles:
\[
 \Pic(\C,\HH)\simeq H^1(\C,\HH^\times).
\]
\end{corollary}
\begin{proof}
A tensor-invertible module over a commutative ring is a finitely generated projective module of rank one. At every stalk $R_c$, it is free by \cref{thm:PID}. Let $\mathcal L$ have tensor inverse $\mathcal M$. Choose mutually dual generators at a stalk and represent them by local sections $s$ and $r$. After shrinking, their pairing is $1$. At every nearby stalk, both modules are free of rank one, and a pairing equal to $1$ makes $s$ a generator. Hence multiplication by $s$ is an isomorphism $\HH\to\mathcal L$ on that neighborhood. Local bases have unit transition functions; changing bases multiplies the cocycle by a coboundary, and gluing copies of $\HH$ constructs the converse. This proves the classification directly.
\end{proof}
The stalk argument is necessary. For a general ringed space with nonlocal stalks, tensor-invertibility need not imply local generation by one section; the distinction is explicitly noted in \cite[Tag 09NT]{Stacks}. We do not suppress that hypothesis by treating our space as locally ringed.

\section{A sharp Hahn Mittag--Leffler theorem}\label{sec:ML}
\subsection{Discrete principal-part data}
Let $D=\{c_j:j\in J\}$ be a closed discrete subset of $\C$, where $J$ is finite or countably infinite. Choose pairwise disjoint ordinary disks $U_j$ centered at $c_j$, and set $U_0=\C\setminus D$. The punctured plane $U_0$ is connected. The sets $U_0,U_j$ cover $\C$; the only nonempty overlaps between distinct members are the punctured disks $U_0\cap U_j$.

Prescribe, for each $j$,
\begin{equation}\label{eq:ppdata}
 p_j(u)=\sum_{\gamma\in S_j}p_{j,\gamma}(u)t^\gamma,
 \qquad p_{j,\gamma}(u)\in u^{-1}\C[u^{-1}],
 \qquad u=z-c_j,
\end{equation}
where $S_j\subseteq\Gamma$ is well ordered and zero principal parts are omitted. The pole order of $p_{j,\gamma}$ is finite but need not be bounded as $j$ or $\gamma$ varies. Each $p_j$ is a legitimate section on its punctured disk.

\begin{definition}[Coherent Mittag--Leffler solution]
A solution for \eqref{eq:ppdata} is a section $B\in\HH(U_0)$ such that $B-p_j$ extends to a section in $\HH(U_j)$ for every $j$.
\end{definition}
The definition imposes extension \emph{coefficientwise across the ordinary center}. It does not pretend that an arbitrary infinite family of meromorphic coefficients already has a finite moving-pole description on the omitted monad.

\begin{theorem}[Uniform-support Mittag--Leffler criterion]\label{thm:ML}
The principal-part data \eqref{eq:ppdata} have a coherent solution if and only if
\begin{equation}\label{eq:MLcriterion}
 S:=\bigcup_{j\in J}S_j\quad\text{is well ordered}.
\end{equation}
When this holds, a solution can be chosen with support contained in $S$. Its coefficients are ordinary meromorphic functions on $\C$ with poles only in $D$. If $S\subseteq\Gamma_{>0}$, the solution and all holomorphic differences $B-p_j$ can be chosen with positive support. Any two solutions differ by a section of $\HH(\C)$.
\end{theorem}
\begin{proof}
Suppose $B=\sum_{\gamma\in T}b_\gamma t^\gamma$ is a solution. If $\gamma\in S_j$ were absent from $T$, the coefficient of $B-p_j$ at $\gamma$ would be the nonzero principal part $-p_{j,\gamma}$. That cannot extend holomorphically to $c_j$. Therefore $S_j\subseteq T$ for every $j$. Since a subset of a well-ordered set is well ordered, \eqref{eq:MLcriterion} is necessary.

Conversely, assume $S$ is well ordered. For each $\gamma\in S$, apply the ordinary Mittag--Leffler theorem to the principal parts $p_{j,\gamma}$, using zero when $\gamma\notin S_j$. Obtain an ordinary meromorphic function $b_\gamma$ on $\C$ with precisely these principal parts and no other poles. These choices are independent at different exponents. Then
\[
 B=\sum_{\gamma\in S}b_\gamma|_{U_0}\,t^\gamma
\]
is a section of $\HH(U_0)$. Every coefficient of $B-p_j$ extends holomorphically to $U_j$, and its support is contained in $S$, so the extensions assemble into a valid section. If all exponents are positive, every section constructed this way has positive support.

Finally, the difference of two solutions extends across each $c_j$. The extensions agree on overlaps, so \cref{prop:sheaf} glues them to a section of $\HH(\C)$. No arbitrary union of the two solutions' local supports is taken in this step.
\end{proof}

\subsection{Why this is not a summability condition over the points}
Condition \eqref{eq:MLcriterion} permits a fixed exponent at infinitely many points. For example, $p_j=t/(z-j)$ is solvable: apply the ordinary Mittag--Leffler theorem once and multiply its result by $t$. The numerical family $(t,t,t,\ldots)$ is not strongly summable, but no sum of that family is required. Ordinary analytic summation across the spatial points takes place \emph{within one coefficient function}.

In the opposite direction, data
\begin{equation}\label{eq:descendingPP}
 p_j(z)=\frac{t^{1/j}}{z-j},\qquad j\ge2,
\end{equation}
are individually single-term Hahn sections, and every finite subproblem is solvable. Their exponent union has a strictly decreasing sequence and is not well ordered. Therefore the infinite problem has no coherent solution.

A common positive lower bound on valuations is insufficient. Replacing $t^{1/j}$ in \eqref{eq:descendingPP} by $t^{1+1/j}$ gives valuations greater than $1$ but still violates \eqref{eq:MLcriterion}. The obstruction concerns order structure throughout the support, not just the smallest magnitude among the data.

\subsection{Classical input and a support-safe implementation}
For completeness, the ordinary construction used in the proof may be arranged as follows. Exhaust $\C$ by closed disks whose boundary circles avoid $D$, group the finitely many new principal parts appearing at each stage, and subtract sufficiently many Taylor terms from each new rational principal part so that its remaining tail is bounded by a prescribed summable positive real sequence on the preceding disk. The resulting ordinary series converges locally uniformly away from $D$ and has the specified principal parts. This is the classical Mittag--Leffler construction \cite{Ahlfors}.

One performs that construction separately for every $\gamma\in S$. The number of subtracted Taylor terms may depend arbitrarily on $\gamma$ and on the spatial point. None of these choices introduces a new Hahn exponent, which is why the construction preserves $S$. There is no requirement of a common real norm estimate for all coefficient functions.

\begin{corollary}[Entire interpolation with exact support criterion]\label{cor:interp}
Let $a_j\in\KK$ be prescribed at a closed discrete set $\{c_j\}$. There exists $A\in\HH(\C)$ with $A(c_j)=a_j$ for all $j$ if and only if $\bigcup_j\supp(a_j)$ is well ordered. When the union is positive, $A$ can be chosen in $\HP(\C)$. The same statement holds for finitely many prescribed derivatives at each point, with the union taken over all prescribed values.
\end{corollary}
\begin{proof}
Necessity follows because evaluation at an ordinary point introduces no new exponent. For sufficiency, interpolate the complex values or finite jets at each exponent by an ordinary entire function, then assemble over the given well-ordered union. Ordinary finite-jet interpolation follows from the classical theorems already used: choose an entire $f_0$ with the desired multiplicity at each point, prescribe the principal part of the desired local jet divided by $f_0$, solve the ordinary Mittag--Leffler problem, and multiply by $f_0$. This creates an entire interpolant with the specified jets.
\end{proof}

\section{The global theorem for moving surcomplex divisors}\label{sec:divisors}
\subsection{Cluster polynomials and the correct compatibility invariant}
Fix a closed discrete ordinary set $D=\{c_j:j\in J\}$ and positive integers $m_j$. At each center prescribe
\begin{equation}\label{eq:clusters}
 P_j(u)=u^{m_j}+\sum_{r=0}^{m_j-1}a_{j,r}u^r,
 \qquad a_{j,r}\in\mm.
\end{equation}
There is no common upper bound on $m_j$. Every root of $P_j$ is infinitesimal. When $\Gamma$ is divisible, prescribing $P_j$ is equivalent to prescribing an unordered finite multiset of elements of $\mm$.

\begin{definition}[Uniform cluster support]
The support of the cluster data is
\begin{equation}\label{eq:SD}
 S_{\mathcal D}=\bigcup_{j\in J}\bigcup_{0\le r<m_j}\supp(a_{j,r})
 \subseteq\Gamma_{>0}.
\end{equation}
The clusters are \emph{uniformly supported} if $S_{\mathcal D}$ is well ordered.
\end{definition}
The definition includes exact zeros $a_{j,r}=0$, which contribute nothing to the union. For finitely many centers, it is automatic, since a finite union of well-ordered sets is well ordered. For infinitely many centers, it is a genuine restriction.

\begin{theorem}[Global moving-divisor criterion]\label{thm:divisor}
For the data \eqref{eq:clusters}, the following are equivalent.
\begin{enumerate}
\item The set $S_{\mathcal D}$ in \eqref{eq:SD} is well ordered.
\item There is a nonzero $F\in\HH(\C)$ such that the complete zero divisor of $F\sh$ on $\C\sh$ consists exactly of the roots of $P_j(z-c_j)$ in the monads of the $c_j$, with their polynomial multiplicities.
\item Given any ordinary entire function $f_0$ whose zero divisor is $\sum_j m_j[c_j]$, there is such an $F$ with coefficient at exponent zero equal to $f_0$, no negative exponents, and support contained in $S_{\mathcal D}^*$.
\end{enumerate}
Equivalently, the prepared polynomial of $F$ at every $c_j$ is exactly $P_j$, and its leading ordinary coefficient has no other zeros. The ordinary function $f_0$ exists by the classical Weierstrass theorem.
\end{theorem}
\begin{proof}
\textit{Necessity.} Suppose $F$ as in (2) exists. Let $\lambda$ be its leading exponent and $f_\lambda$ its leading coefficient. By \cref{cor:roots}, the zeros of $f_\lambda$ are exactly the $c_j$, with orders $m_j$. Put $T=(\supp F-\lambda)\setminus\{0\}$. At every $c_j$, \cref{thm:preparation} constructs a monic polynomial whose lower coefficients have supports in $T^*\setminus\{0\}$. Its roots and their multiplicities are precisely those prescribed by $P_j$, so it equals $P_j$. Therefore
\[
 S_{\mathcal D}\subseteq T^*\setminus\{0\},
\]
a well-ordered set. This proves (2)$\Rightarrow$(1). Multiplying $F$ by an infinite or infinitesimal scalar cannot evade this argument; the normalization by $t^\lambda$ has already removed that freedom.

\textit{Sufficiency.} Suppose (1) holds and fix $f_0$ as in (3). Shrink disjoint disks $U_j$ so that
\[
 q_j(u)=f_0(c_j+u)/u^{m_j}
\]
is an ordinary holomorphic unit on each disk. On its puncture define
\begin{equation}\label{eq:Lj}
 L_j(u)=\LogH\!\left(\frac{P_j(u)}{u^{m_j}}\right)
 =\sum_{n\ge1}\frac{(-1)^{n+1}}n
 \left(\sum_{r<m_j}a_{j,r}u^{r-m_j}\right)^n.
\end{equation}
Let $M=S_{\mathcal D}^*$. Each $L_j$ has support in $M\setminus\{0\}$. At every exponent, its coefficient is a \emph{finite} polynomial in negative powers of $u$. Indeed, finite decomposition in the positive monoid bounds the number of contributing exponent words, and there are only finitely many indices $r$ at each fixed $j$. Thus \eqref{eq:Lj} is precisely a principal-part datum of the kind in \cref{thm:ML}, not an unrestricted Laurent germ.

Apply the positive-support Mittag--Leffler theorem to all the $L_j$. It gives $B\in\HP(U_0)$ with support in $M\setminus\{0\}$ such that
\[
 C_j=B-L_j\in\HP(U_j).
\]
On the covering sets define
\begin{equation}\label{eq:Fglue}
 F_0=f_0\ExpH B\quad\hbox{on }U_0,
 \qquad
 F_j=q_jP_j\ExpH C_j\quad\hbox{on }U_j.
\end{equation}
The exponential law and $\ExpH L_j=P_j/u^{m_j}$ show that these sections agree on every overlap. By \cref{prop:sheaf} they glue to $F\in\HH(\C)$. All expressions have support in $M$ and exponent-zero coefficient $f_0$, so the glued function has these properties too.

On $U_j\sh$, the factors $q_j$ and $\ExpH C_j$ are units. Therefore the zeros and multiplicities are exactly those of $P_j$. On every other monad, $f_0$ is nonzero and hence so is $F\sh$. This proves (1)$\Rightarrow$(3), and (3)$\Rightarrow$(2) is immediate.
\end{proof}

\subsection{Uniqueness modulo global units}
\begin{proposition}[The freedom in realization]\label{prop:realizationtorsor}
Two nonzero sections of $\HH(\C)$ have the same complete zero divisor on $\C\sh$ if and only if their quotient is a global unit. In particular, all realizations in \cref{thm:divisor} differ by multiplication by a unit of $\HH(\C)$. When the exponent-zero leading coefficient is fixed to be the same $f_0$, this unit is uniquely $\ExpH A$ for some $A\in\HP(\C)$.
\end{proposition}
\begin{proof}
One implication follows from the unit criterion. Conversely, apply preparation at each ordinary center. Equality of the zero divisors makes the monic cluster polynomials equal, so the quotient is locally a ratio of unit factors. Away from the leading zeros, both original sections are units after the same kind of local normalization. These local coherent quotients agree on overlaps and glue to a global unit, with an equally coherent inverse. If both leading coefficients are $f_0$ at exponent zero, their ratio has leading term $1$; apply \cref{prop:explog}.
\end{proof}
On $\C$, an ordinary nowhere-zero entire function is $e^h$ for an entire $h$. Hence a general global unit has the form
\[
 t^\lambda e^h\ExpH A,
 \qquad \lambda\in\Gamma,\quad h\in\OO(\C),\quad A\in\HP(\C),
\]
where $h$ is determined up to addition of $2\pi i\Z$. This describes all freedom without choosing an unjustified canonical global product.

\subsection{Root supports versus symmetric supports}
Suppose a cluster has roots $\varepsilon_{j,1},\ldots,\varepsilon_{j,m_j}$. Its coefficients are the elementary symmetric polynomials in these roots, with alternating signs. A well-ordered union of all root supports is sufficient for \eqref{eq:SD}, because finite products use the positive monoid generated by that union. It is not necessary. Symmetric cancellation can remove the descending exponents from the coefficients even when they remain visible in each individual root.

This distinction is intrinsic: the divisor is unordered. A root labeling or a choice of a single branch can destroy precisely the cancellations that make a coherent global equation possible. The next examples show that this is not a marginal phenomenon.

\section{Sharp examples and an explicit nonfactorable divisor splitting}\label{sec:examples}
Throughout this section $t=\omega^{-1}$. The displayed roots lie in $K_\Q$, even when their defining equations have integer Hahn exponents.

\subsection{An impossible sequence of simple zeros}
\begin{example}[Individually valid motions with no exact global realization]\label{ex:simple}
For every integer $j\ge2$, prescribe a single simple zero at
\[
 z_j=j+t^{1/j}.
\]
The cluster polynomial is $P_j(u)=u-t^{1/j}$, so
\[
 S_{\mathcal D}=\{1/2,1/3,1/4,\ldots\}
\]
is not well ordered. By \cref{thm:divisor}, there is no nonzero coherent entire function having \emph{exactly} these simple zeros and no additional zeros on $\C\sh$.
\end{example}
Every finite sub-divisor is realizable, even by a polynomial. The example is therefore a genuinely infinite compatibility obstruction. It does \emph{not} assert that the points $z_j$ cannot all be among the zeros of a coherent entire function. Additional roots can restore the missing symmetric cancellations, as follows.

\subsection{Complete clusters with unbounded multiplicities}
\begin{example}[Uniform equations with nonuniform root scales]\label{ex:fullclusters}
At the center $j\ge2$, prescribe
\begin{equation}\label{eq:unboundedclusters}
 P_j(u)=u^j-t.
\end{equation}
There is only one nonzero lower coefficient, with support $\{1\}$. Thus \cref{thm:divisor} gives a coherent entire function with exactly the clusters
\[
 z_{j,k}=j+\zeta_j^k t^{1/j},\qquad 0\le k<j,
 \qquad \zeta_j=e^{2\pi i/j}.
\]
All roots are simple. Their total multiplicity over the ordinary point $j$ is $j$. The root-support union contains the descending sequence $1/j$, but the symmetric-coefficient support is the singleton $\{1\}$. In fact the entire defining function can be chosen in $\mathscr H_{\Z}(\C)$, although the roots require rational exponents.
\end{example}
This is an exact realization theorem for infinitely many clusters whose sizes are unbounded; it is not a finite-degree polynomial calculation in disguise.

There is a relatively explicit construction in this example. Let
\begin{equation}\label{eq:f0product}
 f_0(z)=\prod_{j=2}^{\infty}
 \left[(1-z/j)\exp\!\left(z/j+\frac{z^2}{2j^2}\right)\right]^j.
\end{equation}
This ordinary Weierstrass product is locally uniformly convergent: on a fixed compact set, the logarithm of its large-$j$ factor is $O(j|z/j|^3)=O(j^{-2})$. Its zeros are exactly $j\ge2$, with multiplicity $j$.

On $U_0=\C\setminus\{2,3,\ldots\}$ put
\begin{equation}\label{eq:bn}
 b_n(z)=-\frac1n\sum_{j=2}^{\infty}\frac1{(z-j)^{jn}},
 \qquad B=\sum_{n\ge1}b_n(z)t^n.
\end{equation}
For each fixed $n$, the ordinary sum over $j$ converges locally uniformly away from the indicated points. For large $j$ on a compact set, $|z-j|\ge j/2$, so its terms are bounded by $(2/j)^{jn}$. The principal part of $b_n$ at $j$ is exactly $-1/[n(z-j)^{jn}]$. Consequently,
\[
 B-\LogH(1-t/(z-j)^j)
\]
has holomorphic coefficients at $j$. The section $f_0\ExpH B$ on $U_0$ extends coherently to all of $\C$ by \eqref{eq:Fglue}, and its local polynomial is $u^j-t$.

Equations \eqref{eq:f0product} and \eqref{eq:bn} involve ordinary convergence inside the coefficient functions. The subsequent sum over $n$ uses Hahn summability. These are two different operations, each justified separately. An unqualified infinite Hahn product over the centers would obscure that distinction.

\subsection{A principal effective divisor with a nonprincipal sub-divisor}
\begin{corollary}[Failure of arbitrary global divisor splitting]\label{cor:splitting}
Let $F$ be a coherent entire function realizing \cref{ex:fullclusters}. Select the root $j+t^{1/j}$ from each cluster. There is no factorization $F=GH$ by coherent entire functions such that $G$ has exactly the selected roots and $H$ has exactly the remaining roots, with the indicated multiplicities. Both prescribed sub-divisors are locally coherent, but neither is globally principal.
\end{corollary}
\begin{proof}
The selected divisor is nonprincipal by \cref{ex:simple}. For the complement, write $a=t^{1/j}$. Its local polynomial is
\[
 Q_j(u)=\frac{u^j-a^j}{u-a}
 =u^{j-1}+a u^{j-2}+\cdots+a^{j-1}.
\]
The coefficient of $u^{j-2}$ is $t^{1/j}$, so its coefficient-support union is not well ordered either. The necessity direction of \cref{thm:divisor} excludes a global realization. Local equations exist on each $U_j$ and are units on punctured overlaps; thus the sub-divisors are locally coherent. A factorization with the prescribed zero divisors would contradict their nonprincipality.
\end{proof}
This does not claim that $F$ is irreducible, that every factorization fails, or that the global ring has a particular unique-factorization property. The assertion concerns this specific, intrinsically meaningful partition of an effective divisor.

\subsection{No repair by enlarging the value group}
The nonexistence in \cref{ex:simple} persists after any set-sized enlargement of $\Gamma$ inside $\No$, and even for a putative section with an arbitrary set-sized surreal support. A larger ambient ordered group does not make the strictly decreasing sequence $1/2,1/3,\ldots$ well ordered. The necessity proof forces that same sequence into one alleged support monoid. In contrast, finite truncation of the \emph{spatial set} removes the obstruction. These two ways of enlarging or approximating the problem must not be confused.

\section{Near-identity automorphisms and interpolation of motions}\label{sec:automorphisms}
For simple zero motions, there is an alternative constructive mechanism. It also illustrates how rich the coherent category remains despite the obstruction theorems.

\begin{theorem}[Global inverse for positive-support perturbations]\label{thm:aut}
Let $U\subseteq\C$ be an ordinary domain and $E\in\HP(U)$. The coherent map
\[
 G(z)=z+E(z)
\]
evaluates to an automorphism of $U\sh$. Its inverse is $H(z)=z+A(z)$ for a unique $A\in\HP(U)$. Both compositions are identities as Hahn-coherent data. If $T=\supp E$, the support of $A$ is contained in $T^*\setminus\{0\}$.
\end{theorem}
\begin{proof}
The standard part of $G(c+\varepsilon)$ is $c$, so evaluation maps $U\sh$ into itself. To construct a right inverse, solve
\begin{equation}\label{eq:inverse-rec}
 A(z)=-E(z+A(z))
 =-\sum_{\gamma\in T}\sum_{n\ge0}
 \frac{e_\gamma^{(n)}(z)}{n!}t^\gamma A(z)^n.
\end{equation}
At $\nu\in T^*\setminus\{0\}$, the terms with $n\ge1$ involve only coefficients of $A$ at exponents smaller than $\nu$, because the external exponent $\gamma$ is strictly positive. The term with $n=0$ is already known. Finite word decomposition makes the coefficient formula finite. Well-founded recursion therefore defines holomorphic coefficients on the same $U$, with the asserted support, and gives $G\circ H=\id$.

For cancellation, suppose two positive perturbations $H_1,H_2$ differ, and let $\nu$ be the least exponent of their difference. In
\[
 G\circ H_1-G\circ H_2
 =(H_1-H_2)+(E\circ H_1-E\circ H_2),
\]
the second summand has no term at $\nu$: every term containing a difference also contains an external positive exponent from $E$. Thus the coefficient at $\nu$ in the displayed difference is nonzero. Composition on the left by $G$ is injective in this class. Associativity of \eqref{eq:comp} gives
\[
 G\circ(H\circ G)=(G\circ H)\circ G=G=G\circ\id,
\]
so $H\circ G=\id$. The same least-exponent argument proves uniqueness among all positive-support inverses, not just those supported in $T^*$.
\end{proof}

\begin{corollary}[Exact interpolation of simple motions]\label{cor:motion}
Let $\{c_j\}$ be closed discrete and prescribe $\delta_j\in\mm$. There is a coherent automorphism $G=\id+E$ with $E\in\HP(\C)$ satisfying $G(c_j)=c_j+\delta_j$ for every $j$ if and only if $\bigcup_j\supp(\delta_j)$ is well ordered.
\end{corollary}
\begin{proof}
Necessity follows by evaluating $E$ at ordinary points. For sufficiency, use \cref{cor:interp} to find $E$ with $E(c_j)=\delta_j$, and apply \cref{thm:aut}.
\end{proof}
If $f_0$ has simple zeros exactly at the $c_j$, then $f_0\circ G^{-1}$ realizes the moved simple divisor. This is consistent with, but weaker than, \cref{thm:divisor}: a bijective change of coordinate cannot split a multiple zero into distinct roots. The moving-divisor theorem permits precisely that splitting and allows arbitrarily large cluster degrees.

For example, $G(z)=z+t z^2$ is a coherent automorphism of the entire finite surcomplex plane, although the same quadratic is not injective on all of $\SC$. Its other algebraic preimage is generally infinite and lies outside $\C\sh$. Its inverse on $\C\sh$ is the positive-support branch
\[
 G^{-1}(z)=z-tz^2+2t^2z^3-5t^3z^4+14t^4z^5-\cdots.
\]
This domain distinction is another reason not to identify coherent-entire functions on the finite plane with all-scale entire functions on the full class field.

\section{Explicit additive and line-bundle obstructions}\label{sec:obstructions}
\subsection{A quotient of sequences of infinitesimals}
Fix the discrete centers $c_j=j$ for $j\ge2$, and the punctured-plane cover from \cref{sec:ML}. Define complex vector spaces
\begin{align}
 \mathcal A_\Gamma&=\prod_{j\ge2}\mm,\label{eq:Aspace}\\
 \mathcal W_\Gamma&=\left\{(a_j)\in\mathcal A_\Gamma:
 \bigcup_{j\ge2}\supp(a_j)\text{ is well ordered}\right\}.
 \label{eq:Wspace}
\end{align}
The second is a vector subspace: finite unions of well-ordered supports are well ordered, and addition can only remove exponents from such a union. No Hahn summation over $j$ is intended in the product \eqref{eq:Aspace}.

An element $a=(a_j)$ determines an additive cocycle on this cover by
\begin{equation}\label{eq:addcocycle}
 p_{0j}(z)=\frac{a_j}{z-j},\qquad p_{j0}=-p_{0j}.
\end{equation}
There are no overlaps among distinct small disks, so no further triple-overlap conditions need to be checked.

\begin{theorem}[A concrete obstruction subspace]\label{thm:obstruction}
The construction \eqref{eq:addcocycle} induces an injective complex-linear map
\begin{equation}\label{eq:injH1}
 \mathcal A_\Gamma/\mathcal W_\Gamma
 \lhook\joinrel\longrightarrow H^1(\C,\HP).
\end{equation}
Exponentiating the transition functions gives an injective group homomorphism
\begin{equation}\label{eq:injPic}
 \mathcal A_\Gamma/\mathcal W_\Gamma
 \lhook\joinrel\longrightarrow\Pic(\C,\HH),
 \qquad
 a\longmapsto\mathcal L(a),
\end{equation}
where $\mathcal L(a)$ is glued by
\begin{equation}\label{eq:transitions}
 g_{0j}(z)=\ExpH\!\left(\frac{a_j}{z-j}\right).
\end{equation}
Its transitions lie in $1+\HP$, so it has an integral model with trivial ordinary reduction. Nevertheless, $\mathcal L(a)$ is nontrivial exactly when $a\notin\mathcal W_\Gamma$.
\end{theorem}
\begin{proof}
An additive coboundary representation on the chosen cover is a system
\[
 p_{0j}=B_0-B_j,
 \qquad B_0\in\HP(U_0),\quad B_j\in\HP(U_j).
\]
This is exactly the positive-support Mittag--Leffler problem for the principal parts $a_j/(z-j)$. By \cref{thm:ML}, it is solvable exactly when the support union is well ordered. Thus the kernel of the map from $\mathcal A_\Gamma$ is $\mathcal W_\Gamma$.

This is a statement about the actual first cohomology class, not just an artifact of one cover. An additive cocycle defines a sheaf torsor. If its class becomes zero after refinement, that torsor has a global section; expressing the global section in the original local trivializations supplies a coboundary on the original cover. Therefore no refinement can remove the obstruction without solving the stated Mittag--Leffler problem.

For clarity, nontriviality after exponentiation can also be checked directly, without appealing to a sheaf cohomology decomposition. Suppose the line bundle with transitions \eqref{eq:transitions} is trivial. There are units $b_0\in\HH(U_0)^\times$ and $b_j\in\HH(U_j)^\times$ satisfying $g_{0j}=b_0/b_j$ on the punctures. Write their leading factors as in \eqref{eq:unitfactor}. Since every $g_{0j}$ has leading term $1$, all leading exponents equal one $\lambda$, and all leading ordinary coefficient functions agree on overlaps. They glue to one nowhere-zero entire $h$. Dividing $b_i$ by $t^\lambda h|_{U_i}$ normalizes every unit to $1+\HP(U_i)$. Taking \(\LogH\) now produces the forbidden additive coboundary unless $a\in\mathcal W_\Gamma$.

Conversely, exponentiating an additive solution trivializes the bundle. The exponential law shows that tensor product corresponds to addition of the sequences, proving \eqref{eq:injPic}. Finally, the same transitions define a locally free module over $\HI$, and their coefficient-zero reductions are all $1$. This is the integral-model statement of \cref{rem:reduction}.
\end{proof}

\begin{example}[An explicit nontrivial bundle over the ordinary plane]\label{ex:bundle}
For $\Gamma=\Q$, take $a_j=t^{1/j}$. Then
\[
 g_{0j}(z)=\sum_{n\ge0}\frac{t^{n/j}}{n!(z-j)^n}
\]
is a valid Hahn unit on every punctured disk. These transitions define a nontrivial line bundle on $(\C,\mathscr H_\Q)$. At each fixed $j$ the support is a perfectly well-ordered arithmetic progression. Nontriviality comes from trying to trivialize all the transitions simultaneously, which would force the descending exponent set $\{1/j:j\ge2\}$ into one global support.
\end{example}
For the bundles in \cref{thm:obstruction}, restriction to any relatively compact ordinary domain is trivial. Only finitely many centers lie in that domain, and the finite principal-part problem is solvable. Principal parts at outside centers are already holomorphic there and are absorbed into the local gauges. Thus a bundle can be trivial on every disk of an ordinary exhaustion while admitting no compatible global trivialization.

\subsection{The divisor obstruction as a line-bundle class}
The cluster data \eqref{eq:clusters} yield unit transitions
\[
 h_{0j}=P_j(z-c_j)/(z-c_j)^{m_j}\in1+\HP(U_0\cap U_j).
\]
Even without a common support across $j$, these transitions define a line bundle. Its logarithmic cocycle is $L_j$ from \eqref{eq:Lj}.

\begin{proposition}[The divisor and logarithmic obstructions coincide]\label{prop:divPic}
The preceding line bundle is trivial if and only if the clusters are uniformly supported. Equivalently, it is trivial exactly when the prescribed effective divisor is globally realizable as in \cref{thm:divisor}.
\end{proposition}
\begin{proof}
The normalization argument in \cref{thm:obstruction} shows that triviality is equivalent to an additive Mittag--Leffler solution for the principal parts $L_j$. By \cref{thm:ML}, this is equivalent to well ordering of $\bigcup_j\supp L_j$. If $S_{\mathcal D}$ is well ordered, all these logarithmic supports lie in $S_{\mathcal D}^*\setminus\{0\}$. Conversely, if the logarithmic union $T$ is well ordered, exponentiating each $L_j$ shows that every coefficient support of $P_j$ lies in $T^*$. Thus $S_{\mathcal D}$ is well ordered. Apply \cref{thm:divisor} for the final equivalence.
\end{proof}
In particular, the effective sub-divisor in \cref{cor:splitting} has a genuine line-bundle obstruction. It is not merely a failure of a particular product formula. The transitions for a moved simple root are $1-a_j/(z-c_j)$, whose logarithm has higher pole terms; they should not be confused with the simpler exponential transitions \eqref{eq:transitions}. Both constructions detect the same well-ordering issue, but their cocycles are not asserted to be equal.

\section{The Picard-group dichotomy}\label{sec:Pic}
\subsection{Separating the ordinary and positive-support factors}
\begin{proposition}[Reduction of the Picard group to positive supports]\label{prop:Picreduction}
There is a natural isomorphism of abelian groups
\begin{equation}\label{eq:PicH1}
 \Pic(\C,\HH)\simeq H^1(\C,\HP).
\end{equation}
In particular the Picard group has a complex vector-space structure, with tensor product corresponding to vector addition. The map \eqref{eq:injPic} is complex-linear with this structure.
\end{proposition}
\begin{proof}
By \cref{cor:Picmeaning,prop:explog},
\[
 \Pic(\C,\HH)
 \simeq H^1(\C,\underline\Gamma)
 \times H^1(\C,\OO^\times)
 \times H^1(\C,\HP).
\]
The first factor is zero: locally constant torsors over the simply connected, locally contractible plane have trivial monodromy and admit a global section. The second is zero by the classical triviality of holomorphic line bundles on $\C$. One may derive the latter from the ordinary exponential exact sequence, $H^1(\C,\OO)=0$, and $H^2(\C,\Z)=0$; see \cite{Demailly}. The remaining factor is cohomology of a complex vector-space sheaf. The exponential construction realizes precisely its inclusion in the product decomposition.
\end{proof}
The conclusion does not say that a contractible space can have a nontrivial ordinary complex line bundle. It says that the coefficient sheaf is different. The obstruction lies in holomorphic dependence coupled to a global well-ordered support, not in ordinary monodromy.

\subsection{Exactly which value groups give vanishing}
\begin{lemma}[Descending positive sequences]\label{lem:descending}
An ordered abelian group $\Gamma$ has no infinite strictly decreasing sequence of positive elements if and only if it is trivial or cyclic.
\end{lemma}
\begin{proof}
For $\Gamma=\Z\delta$ with $\delta>0$, every positive element is a positive integer multiple of $\delta$, so there is no such sequence. The trivial case is immediate.

Conversely, suppose $\Gamma\ne0$. If its positive cone has no least element, recursively choose $\gamma_1>\gamma_2>\cdots>0$. Otherwise let $\delta>0$ be its least positive element. If $\Gamma\ne\Z\delta$, choose a positive $h\notin\Z\delta$. Necessarily $h>n\delta$ for every positive integer $n$: a first integer bound would leave a nonzero remainder strictly between $0$ and $\delta$. Thus $h-\delta,h-2\delta,\ldots$ is a strictly decreasing positive sequence.
\end{proof}

\begin{theorem}[Picard-group vanishing dichotomy]\label{thm:dichotomy}
For every set-sized ordered additive subgroup $\Gamma\subseteq\No$,
\begin{equation}\label{eq:dichotomy}
 \Pic(\C,\HH)=0
 \quad\Longleftrightarrow\quad
 \Gamma=\{0\}\ \text{or}\ \Gamma=\Z\delta\text{ for some }\delta>0.
\end{equation}
\end{theorem}
\begin{proof}
For $\Gamma=0$, the sheaf is $\OO$ and its Picard group vanishes. Suppose $\Gamma=\Z\delta$ with $\delta>0$. A section of $\HP$ is simply
\[
 \sum_{n\ge1}a_n(z)t^{n\delta},\qquad a_n\in\OO(U),
\]
with no condition on the coefficient functions beyond holomorphy: every subset of the positive integer multiples of $\delta$ is well ordered. Let a positive-support additive cocycle be given on an ordinary open cover. Expand it at every $n\delta$. Each coefficient is an ordinary holomorphic additive cocycle. Since $H^1(\C,\OO)=0$, solve its additive Cousin problem on the same cover, obtaining holomorphic local cochains $b_{i,n}$. Their series
\[
 b_i=\sum_{n\ge1}b_{i,n}t^{n\delta}
\]
are valid positive-support sections and solve the original cocycle. The ability to use the same cover follows also from the torsor interpretation: a global section of an ordinary additive torsor gives local coordinates in its original trivializations. Thus $H^1(\C,\HP)=0$, and \cref{prop:Picreduction} proves vanishing.

If $\Gamma$ is noncyclic, choose a strictly decreasing positive sequence $\gamma_j$ by \cref{lem:descending}. The sequence $a_j=t^{\gamma_j}$ belongs to $\mathcal A_\Gamma$ but not to $\mathcal W_\Gamma$. Its line bundle in \cref{thm:obstruction} is nontrivial, proving the other direction.
\end{proof}
The word \emph{cyclic} cannot be replaced by ``has a least positive element''. For example,
\[
 \Gamma=\Z+\Z\omega\subseteq\No
\]
has least positive element $1$, but is not cyclic. The sequence $\gamma_j=\omega-j$ is positive and strictly decreasing. The transitions
\[
 \ExpH\!\left(\frac{t^{\omega-j}}{z-j}\right)
\]
therefore give a nontrivial line bundle. Higher-rank exponents beyond all multiples of the least positive element matter.

\subsection{The size of the nonzero Picard groups}
\begin{theorem}[Continuum many independent classes]\label{thm:dimension}
If $\Gamma$ is noncyclic, then
\begin{equation}\label{eq:lowerdim}
 \dim_\C\Pic(\C,\HH)\ge2^{\aleph_0}.
\end{equation}
If $\Gamma$ is countable and noncyclic, equality holds. In particular,
\begin{equation}\label{eq:Qdim}
 \dim_\C\Pic(\C,\mathscr H_\Q)=2^{\aleph_0}.
\end{equation}
\end{theorem}
\begin{proof}
Choose a decreasing positive sequence $(\gamma_j)$. There is a family $(I_x)_{x\in X}$ of infinite subsets of $\{2,3,\ldots\}$, with $|X|=2^{\aleph_0}$, such that distinct members meet in a finite set. One explicit construction encodes finite binary strings by integers and lets $I_x$ be the set of finite prefixes of an infinite binary sequence $x$. Two distinct infinite sequences share only finitely many prefixes.

For every $x$, define
\[
 a_j^{(x)}=\begin{cases}t^{\gamma_j},&j\in I_x,\\0,&j\notin I_x.\end{cases}
\]
Consider a finite nonzero complex-linear combination of the sequences $a^{(x)}$. Choose one with a nonzero coefficient. Infinitely many indices in its $I_x$ belong to none of the finitely many other selected sets. At each such index the combined sequence still has a nonzero coefficient at $\gamma_j$. These exponents contain an infinite strictly decreasing subsequence. The combination therefore does not belong to $\mathcal W_\Gamma$. The classes of all $a^{(x)}$ are linearly independent in $\mathcal A_\Gamma/\mathcal W_\Gamma$. Now use \cref{thm:obstruction,prop:Picreduction} to obtain \eqref{eq:lowerdim}.

Assume next that $\Gamma$ is countable. On any ordinary open subset, the number of holomorphic functions is at most $\mathfrak c=2^{\aleph_0}$, since a function on each connected component is determined by its values on a countable dense set, and there are at most countably many components. A section of $\HP$ is specified by a subset of a countable group and countably many such functions. Hence $|\HP(U)|\le\mathfrak c$ for every open $U$.

Every torsor can be trivialized on a countable cover selected from a fixed countable basis of the plane. There are at most $\mathfrak c$ such covers. For each one, a cocycle consists of countably many sections, so there are at most $\mathfrak c^{\aleph_0}=\mathfrak c$ cocycles. Thus $|H^1(\C,\HP)|\le\mathfrak c$. Combined with the independent family already constructed, this proves equality of the vector-space dimension with $\mathfrak c$.
\end{proof}
The theorem determines vanishing for every value group and the exact dimension for countable noncyclic groups. It does not claim that \eqref{eq:injH1} is surjective, or provide a canonical basis for all line-bundle classes. Those are different classification questions.

\subsection{Additive cohomology behaves differently}
\begin{proposition}[Nonvanishing for the full Hahn sheaf]\label{prop:fullH1}
If $\Gamma\ne0$, then $H^1(\C,\HH)\ne0$, including when $\Gamma$ is cyclic and the Picard group vanishes.
\end{proposition}
\begin{proof}
Choose $\delta>0$ in $\Gamma$, and prescribe $p_j=t^{-j\delta}/(z-j)$. Each is a section on a punctured disk, but the union $\{-j\delta:j\ge2\}$ is not well ordered. The necessity proof of \cref{thm:ML} shows that its additive cocycle has no global coboundary, hence represents a nonzero class.
\end{proof}
There is no contradiction with the Picard theorem: Hahn exponentiation cannot be applied to these negative-support cochains. Nor is there a contradiction with Cartan's theorem for coherent ordinary analytic sheaves. For $\Gamma\ne0$, the $\OO$-module $\HH$ is not locally finitely generated. At a point $c$, evaluation gives
\[
 R_c/(z-c)R_c\simeq\KK,
\]
which is infinite-dimensional over $\C$; a finitely generated $\OO_c$-module would have a finite-dimensional quotient by its maximal ideal. The adjective ``Hahn-coherent'' specifies the common-support condition, not ordinary module coherence.

\section{Interpretation, boundaries, and further questions}\label{sec:scope}
\subsection{What has been established}
The main result is a local-to-global theorem with an exact obstruction. Ordinary analytic discreteness handles the positions of the standard parts. Well ordering handles the distribution of surreal scales. Neither condition can substitute for the other. The proof succeeds because ordinary Mittag--Leffler summation introduces no new Hahn exponents, while positive-support logarithm converts a nonlinear cluster equation into an additive principal-part problem.

The examples show why this is more than a formal coefficientwise reformulation. An individually infinitesimal perturbation at each point can be globally incompatible even with a common positive lower bound on all valuations. Conversely, roots with incompatible individually displayed scales can belong to a compatible global equation when their symmetric coefficients cancel. The line-bundle theorem then turns this support obstruction into a global invariant, and the stalk theorem shows that it survives despite particularly simple local algebra.

\begin{center}\small
\begin{tabular}{@{}L{0.26\textwidth}L{0.32\textwidth}L{0.32\textwidth}@{}}
\toprule
Question & Exact positive result & Sharp boundary\\
\midrule
Principal parts at a discrete set & Realizable exactly for a well-ordered union of principal-part supports & A descending sequence of exponents prevents a coherent solution.\\[0.45em]
Finite moving clusters & Realizable exactly for a well-ordered union of symmetric coefficient supports & Individually well-defined roots need not form a principal divisor.\\[0.45em]
Ordinary-base local algebra & Every stalk is a principal ideal domain & It is nonlocal for every nonzero value group.\\[0.45em]
Simple infinitesimal motions & Uniform support gives a near-identity coherent automorphism & A bijection cannot split a multiple zero.\\[0.45em]
Line bundles over the plane & Picard vanishing is equivalent to a trivial or cyclic value group & Every noncyclic group gives continuum many independent classes.\\[0.45em]
Ordinary reduction & Positive-unit transition bundles have trivial integral reduction & Trivial reduction does not imply triviality over $\HH$.\\
\bottomrule
\end{tabular}
\end{center}

\subsection{What the conclusions do not say}
All cohomology groups and line bundles here live over the ordinary complex plane with its usual topology. We have not assigned ordinary sheaf cohomology to the proper class of all surcomplex points. Each fixed coefficient group is a set. Every algebraic support and recursive construction is set-sized.

The theorem for global divisors concerns $\C\sh$, the finite surcomplex plane. It neither contradicts nor weakens the all-scale polynomial rigidity theorem in \cite[\S12]{Source}. A function constructed by \cref{thm:divisor} has a coherent description at the ordinary scale; it need not have coherent descriptions at every infinite surreal scale.

A Hahn principal-part solution on $\C\setminus D$ is not automatically a meromorphic function on every omitted monad with only finitely many poles. For moving divisors, finite polynomial cancellation supplies that extra structure. For the more general Mittag--Leffler theorem, no such additional conclusion is claimed.

Nor does the article claim to solve the questions about canonical global exponentiation, sectorial summation, or essential singularities in other surreal frameworks. Its selected research problem is the global compatibility of the coefficient sheaf specified in the attachment. The exact solutions and obstructions proved here are independent of those broader choices.

\subsection{Further precise directions}
The explicit injection $\mathcal A_\Gamma/\mathcal W_\Gamma\hookrightarrow H^1(\C,\HP)$ leaves a natural structural problem: characterize the cokernel, or find a useful normal form for an arbitrary class. Our proof of the Picard-group dimension for countable groups counts classes but does not supply such a normal form.

For uncountable value groups, the lower bound $2^{\aleph_0}$ need not be optimal. Determining the dimension in terms of cardinal and order invariants of $\Gamma$ would refine the dichotomy. The positive cone, not just the cardinality or the existence of a least positive element, should remain relevant.

There is also a multivariable extension problem: local Weierstrass division with parameter coefficients can be combined with finite-module algebra, but an exact global support criterion for higher-codimension analytic sets requires more than the one-variable logarithmic reduction. The present article provides a test case and explicit obstructions that any such extension should preserve; it does not claim that result here.

\section{Conclusion}
The globally coherent category has a precise global compatibility boundary. Infinitely many infinitesimal data are compatible precisely when the appropriate coefficient supports can coexist inside one well-ordered set. For divisors, the relevant coefficients are the symmetric coefficients of the local cluster polynomials. For additive principal parts, they are the principal-part exponents themselves. For line bundles, positive logarithms reveal the same issue as a first-cohomology class.

This yields a global Weierstrass theorem with an exact converse, explicit principal divisors with nonprincipal sub-divisors, and a complete vanishing dichotomy for the Picard group of the coefficient sheaf. The phenomena already occur in the countable divisible workspace $\Gamma=\Q$. They require neither proper-class summation nor a problematic notion of fine-topological convergence.

\appendix
\section{Support and dependency audit}\label{app:audit}
The following bookkeeping summarizes the operations for which infinite formal expressions occur. A support in the right column is a containing support; cancellations may make the actual support smaller.

\begin{center}\small
\begin{longtable}{@{}L{0.36\textwidth}L{0.53\textwidth}@{}}
\toprule
Construction & Controlling support and finiteness mechanism\\
\midrule
\endhead
Evaluation at $c+\varepsilon$ & $S+(\supp\varepsilon)^*$; finite exponent decompositions justify Taylor regrouping.\\[0.45em]
Positive exponential or logarithm & $S^*$ for positive well-ordered $S$; finitely many exponent words contribute at each exponent.\\[0.45em]
Preparation at any ordinary center & $T^*$, where $T$ is the relative positive support of the original section; the same monoid works at all centers.\\[0.45em]
Mittag--Leffler solution & The union of the principal-part supports; ordinary analytic choices at a fixed exponent introduce no new exponents.\\[0.45em]
Logarithm of a cluster polynomial & $S_{\mathcal D}^*$; the polynomial has finitely many coefficients at each center, so every Hahn coefficient remains a finite principal part.\\[0.45em]
Entire divisor realization & $S_{\mathcal D}^*$ throughout the additive solution, exponentiation, and gluing.\\[0.45em]
Near-identity inverse & $(\supp E)^*$; external positive exponents make the coefficient recursion triangular.\\[0.45em]
Cyclic-group additive coboundary & The universal set $\{\delta,2\delta,\ldots\}$; every coefficientwise solution is automatically admissible.\\[0.45em]
Nontriviality arguments & A forced descending exponent sequence would have to be a subset of one legitimate support, which is impossible.\\
\bottomrule
\end{longtable}
\end{center}

The dependency chain is
\[
 \begin{gathered}
 \text{Hahn support calculus}\ \Longrightarrow\
 \text{units, logarithms, evaluation, preparation},\\
 \text{ordinary Mittag--Leffler}+\text{support control}\ \Longrightarrow\
 \text{the sharp Hahn Mittag--Leffler criterion},\\
 \text{preparation}+\text{logarithmic principal parts}\ \Longrightarrow\
 \text{the moving-divisor criterion},\\
 \text{preparation}\ \Longrightarrow\ \text{principal-ideal-domain stalks},\\
 \text{unit splitting}+\text{Cousin theory}+\text{descending supports}\
 \Longrightarrow\ \text{the Picard-group dichotomy}.
 \end{gathered}
\]
The principal-ideal-domain argument uses the Euclidean algorithm on \emph{finite-degree} polynomials and a minimum of a nonempty set of nonnegative integers. It does not invoke an unproved Noetherian property of an infinite coefficient ring. The Picard argument does not silently assume that the stalks are local; local rank-one triviality was established separately.

\section{Reproducible symbolic checks}\label{app:checks}
The accompanying \texttt{verify\_examples.py} uses exact symbolic algebra to check finite truncations and polynomial identities that illustrate the proofs. It checks a mixed principal-part logarithm and exponential through degree six in $t$, cancellation for two finite root clusters, the two-sided inverse of $z+t z^2$ through degree seven, complementary root-cluster factors for degrees two through ten, finite Lagrange interpolation of positive-support motions, and a truncated instance of the preparation recursion. The output is recorded in \texttt{verification\_report.txt}.

For the preparation check, take
\[
 F(u)=u^2+t(1+u^3)+t^2(2u+u^4).
\]
The script constructs $A_n$ of degree less than two and polynomial $B_n$ from
\[
 h_n=F_n-\sum_{i=1}^{n-1}A_iB_{n-i},\qquad
 A_n=R_2h_n,\quad B_n=D_2h_n,
\]
and checks $(u^2+A)(1+B)=F$ through degree six. This is a finite example of the well-founded recursion in \cref{thm:preparation}.

For exponential coefficients, the script uses the exact recurrence
\[
 e_0=1,\qquad
 n e_n=\sum_{k=1}^{n}k b_k e_{n-k},
 \qquad \ExpH\left(\sum_{k\ge1}b_kt^k\right)=\sum_{n\ge0}e_nt^n.
\]
This avoids numerical error and makes every truncation order explicit. No finite symbolic computation establishes the nonexistence of a well ordering for an infinite descending sequence, the Picard-group dimension, or the validity of an arbitrary-support recursion. Those assertions are proved in the main text. The script is a consistency check on formulas, not a proof assistant or a novelty checker.

\section{Notation and provenance}
\begin{center}\small
\begin{tabular}{@{}L{0.26\textwidth}L{0.63\textwidth}@{}}
\toprule
Symbol & Meaning\\
\midrule
$\SC=\No[i]$ & Algebraically closed surcomplex class field.\\[0.3em]
$t^\gamma$ & Conway monomial $\omega^{-\gamma}$; not an independently chosen exponential.\\[0.3em]
$\KK$, $\mm$ & Fixed complex Hahn field and its positive-valuation infinitesimals.\\[0.3em]
$\HH$ & Sheaf of holomorphic Hahn coefficient data on the ordinary base.\\[0.3em]
$\HP$, $\HI$ & Positive-support additive sheaf and nonnegative-support subring sheaf.\\[0.3em]
$U\sh$, $F\sh$ & Infinitesimal thickening of an ordinary domain and Taylor--Hahn evaluation.\\[0.3em]
$S^*$ & Additive monoid generated by a positive well-ordered set $S$, including zero.\\[0.3em]
$P_j$ & Monic polynomial encoding the complete infinitesimal root cluster over $c_j$.\\[0.3em]
$R_c$ & Stalk of $\HH$ at an ordinary base point, not a fine germ at one surcomplex point.\\[0.3em]
$\mathcal A_\Gamma/\mathcal W_\Gamma$ & Explicit quotient of infinitesimal sequences by uniformly supported sequences.\\[0.3em]
$\Pic(\C,\HH)$ & Tensor-invertible sheaves on the indicated ringed space; equivalently rank-one locally free sheaves by \cref{cor:Picmeaning}.\\
\bottomrule
\end{tabular}
\end{center}

The source manuscript \cite{Source} supplied the coherent framework, evaluation, support-preserving gluing, preparation, and exact local root counts. Those inputs have been identified at their points of use. The global compatibility criteria and cohomological conclusions are the developments made in this continuation. The article is written as a self-contained research draft; external references are used for established foundational and ordinary analytic results, not as certificates of the new claims.

\begin{thebibliography}{99}
\addcontentsline{toc}{section}{References}

\bibitem{Source}
\emph{Surcomplex Analysis: Infinitesimal Calculus, Hahn-Coherent Holomorphy, and Contour Theory}.
User-supplied manuscript, September 21, 2026, file
\texttt{surcomplex\_analysis(3).tex}.
The present article uses its \S5 (coherent functions and gluing), \S7 (preparation and root lifting), and the distinction between finite-scale and all-scale coherence. Authorship and publication status were not supplied.

\bibitem{Ahlfors}
Lars V. Ahlfors,
\emph{Complex Analysis}, 3rd edition,
McGraw--Hill, 1979.
Classical Weierstrass products, Mittag--Leffler theory, and elementary holomorphic function theory.

\bibitem{BM}
Alessandro Berarducci and Vincenzo Mantova,
``Transseries as germs of surreal functions,''
\emph{Transactions of the American Mathematical Society} \textbf{371} (2019), 3549--3592.
\href{https://arxiv.org/abs/1703.01995}{arXiv:1703.01995}.

\bibitem{CE}
Ovidiu Costin and Philip Ehrlich,
``Integration on the surreals,''
\emph{Advances in Mathematics} \textbf{452} (2024), article 109823.
\href{https://arxiv.org/abs/2208.14331}{arXiv:2208.14331}.

\bibitem{Demailly}
Jean-Pierre Demailly,
\emph{Complex Analytic and Differential Geometry}, online book.
In particular, Chapter V, \S11 (holomorphic vector bundles), and Chapter IX (cohomological vanishing on Stein spaces).
\href{https://www-fourier.univ-grenoble-alpes.fr/~demailly/manuscripts/agbook.pdf}{Author's full text}, accessed September 21, 2026.

\bibitem{EK}
Philip Ehrlich and Elliot Kaplan,
``Surreal ordered exponential fields,''
\emph{The Journal of Symbolic Logic} \textbf{86} (2021), no.~3, 1066--1115.
\href{https://doi.org/10.1017/jsl.2021.59}{doi:10.1017/jsl.2021.59};
\href{https://arxiv.org/abs/2002.07739}{arXiv:2002.07739}.
The surcomplex exponentiation discussion is in \S11 of the arXiv version.

\bibitem{MS}
J\"orn M\"uller and Alexander Strohmaier,
``The theory of Hahn-meromorphic functions, a holomorphic Fredholm theorem, and its applications,''
\emph{Analysis \& PDE} \textbf{7} (2014), no.~3, 745--770.
\href{https://doi.org/10.2140/apde.2014.7.745}{doi:10.2140/apde.2014.7.745};
\href{https://arxiv.org/abs/1205.0236}{arXiv:1205.0236}.

\bibitem{Neumann}
B. H. Neumann,
``On ordered division rings,''
\emph{Transactions of the American Mathematical Society} \textbf{66} (1949), 202--252.

\bibitem{Poonen}
Bjorn Poonen,
``Maximally complete fields,''
\emph{L'Enseignement Math\'ematique} (2) \textbf{39} (1993), 87--106.
\href{https://math.mit.edu/~poonen/papers/amsval.pdf}{Author's full text}.
Corollary~4 supplies the algebraic-closure statement for divisible Hahn value groups.

\bibitem{Stacks}
The Stacks Project Authors,
\emph{The Stacks Project}, Section 20.6, ``First cohomology and invertible sheaves,''
\href{https://stacks.math.columbia.edu/tag/09NT}{Tag 09NT}, accessed September 21, 2026.
The local-stalk hypothesis and the distinction between tensor-invertibility and rank-one local freeness on arbitrary ringed spaces are relevant here.

\bibitem{vdDE}
Lou van den Dries and Philip Ehrlich,
``Fields of surreal numbers and exponentiation,''
\emph{Fundamenta Mathematicae} \textbf{167} (2001), 173--188.
\href{https://doi.org/10.4064/fm167-2-3}{doi:10.4064/fm167-2-3}.
Erratum, \emph{Fundamenta Mathematicae} \textbf{168} (2001), 295--297.

\end{thebibliography}
\end{document}
